% !TeX root = ../教材题库_全部条目_转写工作稿.tex
% Source: XJ-B1-C04-HW-P136-Y0953
% 第4章 / 4.4 函数与方程 / 教材 p.136 / PDF 第143页
% OCR未校对：公式与文字均可能有误；未完成数学TeX转写；禁止直接发给学生
% 原题及相邻原解截图（核对用）：../题目原图/XJ-B1-C04-HW-P136-Y0953-01.webp
% 本文件只在转写完成后作为题干使用；OCR 条目尚不是可靠题干。
\begin{bankitem}{习题4.4 第5题}{第4章 / 4.4 函数与方程 / 教材 p.136 / PDF 第143页}{XJ-B1-C04-HW-P136-Y0953}{OCR未校对：公式与文字均可能有误；未完成数学TeX转写；禁止直接发给学生}
\begin{ocrdraft}
[教材p.136 / PDF第143页]
\par
5．若<b<0则函数丆(x）一（一）(x一仞十（工一仞（一c）十(x一园(x一
\par
的两个零点分别位于区间（
\par
(A)（山和仂，c)内
\par
(C)仂，0和（c，十亠）内
\par
(D)（一亠，和（0十亠）内
\end{ocrdraft}
\end{bankitem}
